\documentclass[11pt,a4paper,answers]{exam}
\usepackage[T1]{fontenc}
\usepackage{longtable}
\usepackage{graphicx}

\usepackage{hyperref}
%\urlstyle{same}  % don't use monospace font for urls
\setlength{\parindent}{0pt}
\setlength{\parskip}{6pt plus 2pt minus 1pt}

\usepackage{geometry}
\usepackage{multicol}
\usepackage{relsize}
\usepackage{pgf}
\usepackage{fancyvrb}

\geometry{
  %includeheadfoot,
  margin=2.54cm
}
\title{Worksheet 02: Further exercises on Java Reflection}
\author{}
\date{}

\fvset{tabsize=2}
\fvset{fontsize=\relsize{-1}}
\fvset{frame=single}

\begin{document}

\maketitle

The Reflection API allows a Java program to inspect and manipulate itself; it comprises the \verb!java.lang.Class! 
class and the \verb!java.lang.reflect! package, which represents the members of a class with \verb!Method!, 
\verb!Constructor!, and \verb!Field! objects.

\begin{questions}
\question
Write a Java program that reads the name of a class from the command line and emits the interface of the class in Java 
syntax (interface or class, modifiers, constructors, methods, fields; no method bodies). 

Hint: You can load a class whose name you know with \verb!java.lang.Class.forName()!. The 
\verb!java.lang.Class! class offers a rich interface that enables you to inspect the interface of any class.

Apply this program to a set of classes and interfaces as test input. You may also use the program on itself.

\question
Write a program that reads a class name and a list of arguments, 
and creates an object of that class where the read arguments are passed to the constructor.

Hint: Treat arguments as strings. A \verb!java.lang.Class! can enumerate its constructors. 
Choose a constructor with the appropriate parameter count. 
Then, find the parameter types. To create typed argument objects, call the appropriate constructors 
that take a string as their only argument. 
Call dynamic constructors using 

\verb!java.lang.reflect.Constructor.newInstance()!.

\question
Write a \verb!JUnit! test to help grade the internal correctness of a student's submitted program for a hypothetical 
assignment.

Make the tests fail if the class under test has any of the following:
\begin{itemize}
\item more than four fields
\item any non-private fields
\item any fields of type \verb!ArrayList!
\item fewer than two private helper methods
\item any method that has a \verb!throws! clause
\item any method that returns an \verb!int!
\item missing a zero-argument constructor	
\end{itemize}

\question
Normally it is up to the programmer to write a \verb!toString()! method for each class one creates. 
This exercise is about writing a general \verb!toString! method once and for all. 
As part of the reflection API for java, it is possible to find out which fields exist for a given object, 
and to get their values. The purpose of this exercise is to make a \verb!toString! method that gives a printed 
representation of any object, in such a manner that all fields are printed, 
and references to other objects are handled as well. 

To solve this exercise, you will need to examine the \verb!java.lang.reflect! API. 
You will (probably) end up with around 50--60 lines of code, including that necessary for trying it out. 
\end{questions}

\end{document}